\notechapter{N-20220501}{Lagrange Multipliers, Hessians, Bordered Hessians, and Degenerate Extrema}

\section{Why this note was too compressed before}

The note is not just “Lagrange multipliers plus Hessian.”  It is a long reflection on how to decide extrema of multivariable functions.  It contains the method, the derivation, examples, doubts, bordered Hessian, degeneracy, higher-order tests, numerical methods, MATLAB, and links to optimization theory.  The repaired version preserves that exploration.

\section{Lagrange multipliers: the basic idea}

Suppose we want to find extrema of
\[
  f(x_1,\ldots,x_n)
\]
under the constraint
\[
  g(x_1,\ldots,x_n)=0.
\]
The Lagrange function is
\[
  L(x_1,\ldots,x_n,\lambda)
  =f(x_1,\ldots,x_n)-\lambda g(x_1,\ldots,x_n).
\]
The stationarity equations are
\[
  \frac{\partial L}{\partial x_i}=0,
  \qquad i=1,\ldots,n,
  \qquad
  \frac{\partial L}{\partial \lambda}=0.
\]
Equivalently,
\[
  \nabla f=\lambda \nabla g,
  \qquad
  g=0.
\]
Geometrically, at a constrained extremum, the level surface of $f$ is tangent to the constraint surface.  Therefore their normal directions are parallel.  This explains why gradients are proportional.

The original Chinese note says that this method is both surprising and confusing: surprising because it transforms a constrained problem into a stationarity system, confusing because the multiplier seems to appear out of nowhere.  The geometric explanation removes much of the mystery: $\lambda$ is the scalar measuring how the two normal vectors compare.

\section{Example: minimizing distance to a line}

Consider
\[
  f(x,y)=x^2+y^2
\]
under the constraint
\[
  x+y=1.
\]
Set
\[
  L(x,y,\lambda)=x^2+y^2-\lambda(x+y-1).
\]
The equations are
\[
  2x-\lambda=0,
  \qquad
  2y-\lambda=0,
  \qquad
  x+y=1.
\]
The first two give $x=y$, and therefore $x=y=1/2$.  This is the closest point of the line $x+y=1$ to the origin.

The note remarks that this example is simple but educational.  For nonlinear constraints, the stationarity system can become hard to solve; in practice one may need numerical methods such as Newton iteration.  The Lagrange method gives equations, not automatic closed forms.

\section{Why the method is valid}

A more rigorous proof uses the implicit function theorem.  If $\nabla g(p)\ne0$, then near $p$ the constraint $g=0$ is a smooth hypersurface.  A tangent vector $v$ to the constraint satisfies
\[
  \nabla g(p)\cdot v=0.
\]
If $p$ is a constrained extremum of $f$, then the derivative of $f$ along every tangent direction must vanish:
\[
  \nabla f(p)\cdot v=0
  \qquad \text{for all tangent }v.
\]
Thus $\nabla f(p)$ is perpendicular to the tangent space.  The normal space is spanned by $\nabla g(p)$, so
\[
  \nabla f(p)=\lambda \nabla g(p).
\]
For several constraints $g_1=\cdots=g_m=0$, the condition becomes
\[
  \nabla f=\sum_{i=1}^m \lambda_i \nabla g_i,
\]
assuming the constraint gradients are independent.

This is where the note's later warning about singular constraints enters: if the gradients of constraints are linearly dependent, the usual bordered-Hessian test may fail.

\section{Hessian matrix and the second derivative test}

For a twice differentiable function $f(x_1,\ldots,x_n)$, the Hessian matrix at $x_0$ is
\[
  H_f(x_0)=\left(\frac{\partial^2 f}{\partial x_i\partial x_j}(x_0)\right)_{i,j}.
\]
If the mixed partials are continuous, the Hessian is symmetric.  This lets us use the linear algebra of symmetric matrices.

At an unconstrained critical point $\nabla f(x_0)=0$:
\begin{enumerate}[(i)]
  \item if $H_f(x_0)$ is positive definite, $x_0$ is a local minimum;
  \item if $H_f(x_0)$ is negative definite, $x_0$ is a local maximum;
  \item if $H_f(x_0)$ is indefinite, $x_0$ is a saddle point;
  \item if $H_f(x_0)$ is singular or semidefinite, the second derivative test is inconclusive.
\end{enumerate}

The proof comes from Taylor expansion:
\[
  f(x_0+h)=f(x_0)+\frac12 h^T H_f(x_0)h+o(\|h\|^2)
\]
when the gradient term vanishes.  Thus the sign of the quadratic form controls the local shape.

\section{Sylvester criterion and eigenvalues}

For a real symmetric matrix, positive definiteness can be tested by eigenvalues or by leading principal minors.  Sylvester's criterion says that $A$ is positive definite iff all leading principal minors are positive.  Negative definiteness has alternating signs.  Eigenvalues give an equivalent test: positive definite means all eigenvalues are positive; indefinite means there are both positive and negative eigenvalues.

The note explicitly connects this to linear algebra.  A Hessian is not just an array of second derivatives; it is a quadratic form.  The eigenvectors are principal directions of curvature, and the eigenvalues measure curvature along those directions.

\section{Bordered Hessian}

For constrained extrema, the ordinary Hessian of $f$ alone is not enough because we are only allowed to move along directions tangent to the constraint.  The note introduces the bordered Hessian as a way to add constraint-gradient information to the second derivative test.

For constraints $g_i(x)=0$, define the Lagrangian
\[
  L=f-\sum_{i=1}^m \lambda_i g_i.
\]
A common bordered Hessian has the block form
\[
  B=
  \begin{pmatrix}
  0_{m\times m} & (Dg)^T\\
  Dg & H_x L
  \end{pmatrix},
\]
where $Dg$ records gradients of the constraints and $H_xL$ is the Hessian with respect to the original variables.  Different texts use slightly different sign conventions and ordering of rows, so one must check the convention before applying a sign rule.

The intuitive meaning is: the border records which directions are forbidden by the constraints, and the Hessian block records curvature of the Lagrangian in the remaining directions.

\section{Example revisited with a bordered Hessian}

For the simple constraint $x+y=1$, the constraint gradient is
\[
  \nabla g=(1,1).
\]
If $f(x,y)=x^2+y^2$, then
\[
  H_f=
  \begin{pmatrix}
  2&0\\
  0&2
  \end{pmatrix}.
\]
The tangent direction to $x+y=1$ is spanned by $(1,-1)$.  Along that direction,
\[
  (1,-1)H_f(1,-1)^T=4>0,
\]
so the constrained critical point is a constrained minimum.

This tangent-space calculation is often clearer than memorizing signs of bordered determinants.  The bordered Hessian is a compact determinant method for doing the same restriction systematically.

\section{Degenerate cases and higher-order tests}

The note explicitly worries about cases where the Hessian is singular or the bordered Hessian test is inconclusive.  In such cases, second-order information is not enough.  One must examine higher-order Taylor terms:
\[
  f(x_0+h)=f(x_0)+P_k(h)+o(\|h\|^k),
\]
where $P_k$ is the first nonzero homogeneous term after the constant term.  If $P_k$ has a fixed sign, one may still obtain an extremum; if it changes sign, one gets a saddle-type point.

A standard example is $f(x,y)=x^4+y^4$, whose Hessian at the origin is zero, but the origin is still a strict local minimum.  Conversely, $f(x,y)=x^4-y^4$ also has zero Hessian at the origin but has a saddle.  Thus degeneracy forces us to look beyond quadratic approximation.

\section{Saddle points and geometric visualization}

A saddle point is a critical point that is neither a local maximum nor a local minimum.  If the Hessian has both positive and negative eigenvalues, then some directions curve upward and others downward.  The note suggests using contour plots to detect saddle behavior.  This is a good habit: level curves reveal local shape far more clearly than raw formulas.

The classical picture is $f(x,y)=x^2-y^2$.  Along the $x$-axis it increases; along the $y$-axis it decreases.  The origin is therefore a saddle.

\section{Optimization beyond equality constraints}

The note mentions Kuhn--Tucker conditions, penalty methods, and interior-point methods.  These belong to constrained optimization with inequalities.  If constraints include
\[
  h_j(x)\le0,
\]
then Lagrange multipliers are supplemented by nonnegativity and complementary slackness:
\[
  \mu_j\ge0,
  \qquad
  \mu_jh_j(x)=0.
\]
This says that an inequality constraint contributes a force only when it is active.  This is the natural extension of the Lagrange multiplier idea.

Penalty methods instead replace constraints by large penalty terms, for example
\[
  f(x)+\rho \sum_j h_j(x)^2.
\]
Interior-point methods keep the iterate inside the feasible region using barrier functions.  All of these methods share the same philosophy: convert a constrained problem into a system that can be treated with unconstrained or nearly unconstrained tools.

\section{Global topology and extrema}

The note speculates about global topology.  This is not random.  Local Hessian information tells us what happens near a critical point, but global existence of maxima and minima depends on compactness.  A continuous function on a compact set attains maximum and minimum.  Thus topology enters before calculus even begins.

There are also deeper links: Morse theory studies how critical points of smooth functions reflect the topology of a manifold.  The Hessian at a nondegenerate critical point gives the Morse index, and the collection of critical points helps reconstruct the space.  This is the high-level version of the note's intuition that extrema are connected with global structure.

\section{Computation and MATLAB}

The note ends with computational reflection.  For high-dimensional functions, one can use symbolic or numerical software to compute gradients, Hessians, eigenvalues, and bordered Hessians.  A typical workflow is:
\begin{enumerate}[(i)]
  \item define $f$ and constraints $g_i$;
  \item solve $\nabla f=\sum \lambda_i\nabla g_i$ together with $g_i=0$;
  \item compute the Hessian of the Lagrangian;
  \item restrict to the tangent space or use a bordered Hessian;
  \item inspect eigenvalues or determinant signs.
\end{enumerate}
The note's MATLAB idea is therefore mathematically meaningful: computation is not replacing theory, but automating the linear algebra after the theoretical conditions have been derived.


\section{Second variation on the tangent space}

For constrained extrema, the cleanest conceptual test is to restrict the Hessian of the Lagrangian to tangent directions.  If the constraint is \(g(x)=0\), then tangent vectors \(v\) satisfy
\[
  \nabla g(x_0)\cdot v=0.
\]
The second variation is
\[
  v^T H_xL(x_0,\lambda)v.
\]
If this quadratic form is positive definite on the tangent space, the point is a constrained local minimum; if negative definite, a constrained local maximum.

This explains what the bordered Hessian is doing: it is a determinant device for testing the same restricted quadratic form.

\section{A useful afterword}

This note is about the transition from calculus to optimization and differential geometry.  Lagrange multipliers are normal-vector equations.  Hessians are quadratic forms.  Bordered Hessians restrict second variation to tangent directions.  Degenerate cases require higher-order terms.  Inequality constraints lead to KKT theory.  Compactness guarantees global extrema, and Morse theory relates critical points to topology.  The original confusion is productive: it points exactly toward the deeper structure behind the computational rules.


\paragraph{Expanded synthesis.}
For this chapter, the elementary material should be read as a study of what remains invariant when coordinates change. The earlier sections (`Why this note was too compressed before`, `Lagrange multipliers: the basic idea`, `Example: minimizing distance to a line`, `Why the method is valid`, `Hessian matrix and the second derivative test`) compute with arrays, bases, pivots, determinants, or eigenvectors; the deeper object is a linear map together with the spaces it acts on. A matrix is a representation of that map after bases have been chosen. This is why formulas such as
\[
  A' = P^{-1}AP,\qquad \widetilde A = T^{-1}AS
\]
are not cosmetic: they distinguish an intrinsic morphism from a coordinate description.

The structural hierarchy is as follows. Kernels record what a map forgets, images record what it reaches, cokernels record obstructions to solvability, and quotients record information after an equivalence has been imposed. Determinants act on the top exterior power and therefore measure oriented volume; traces are invariant under cyclic rearrangement and become characters in representation theory; adjoints depend on an inner product and lead to self-adjoint spectral theory.

A useful way to return to the computations is to ask, for every formula, which invariant it is measuring. Row reduction measures rank and kernel; diagonalization measures decomposition into invariant subspaces; Gram--Schmidt measures orthogonal projection; positive definiteness measures ellipsoid geometry; tensor notation measures how components transform. The elementary methods are therefore the finite-dimensional laboratory in which duality, quotienting, functoriality, and spectral decomposition can be seen clearly.
